\chapter*{ABSTRAK}
\addcontentsline{toc}{chapter}{ABSTRAK}

\begin{center}

    \textbf{Pengembangan dan Implementasi Antarmuka Pengguna Sistem Pelacakan ITB Ultra-Marathon} \\
    \vspace{0.4cm}
    oleh \\
    \vspace{0.4cm}
    Dinda Thalia Fahira\\
    18222055\\
\end{center}


\begin{spacing}{1.0}
Penelitian ini bertujuan untuk mengembangkan dan mengimplementasikan antarmuka pengguna (\textit{front-end}) sistem pelacakan ITB Ultra-Marathon yang bernama UltraTrack. Sistem ini dirancang secara khusus untuk memvisualisasikan posisi dan pergerakan peserta secara \textit{real-time}, serta menyediakan antarmuka pemantauan komprehensif yang disesuaikan untuk berbagai kelompok pengguna, termasuk peserta, ketua tim, panitia, \textit{supporter}, dan penonton. Pengembangan sistem ini dilatarbelakangi oleh tingginya kompleksitas operasional pada ajang \textit{ultra-marathon} berskala besar di Indonesia. Sistem pelacakan yang digunakan pada penyelenggaraan sebelumnya masih bersifat pasif, terfragmentasi, dan mengandalkan proses manual. Keterbatasan tersebut mengakibatkan panitia kesulitan dalam memantau posisi ribuan peserta secara serentak, mengoordinasikan logistik tim, serta merespons situasi darurat dengan cepat di sepanjang rute yang menantang.

Proses pengembangan dilakukan menggunakan pendekatan \textit{Software Development Life Cycle} (SDLC) model \textit{Waterfall}. Berdasarkan hasil pemilihan teknologi menggunakan metode \textit{Analytic Hierarchy Process} (AHP) dengan skor konsistensi (CR) 0,016, sistem ini dibangun secara \textit{cross-platform}. Aplikasi \textit{mobile} berbasis React Native/Expo dikembangkan untuk pelacakan peserta di lapangan, sedangkan aplikasi web berbasis React/Vite dibangun untuk pemantauan terpusat. Perancangan antarmukanya secara khusus mengadopsi prinsip \textit{night-optimized} dan \textit{glanceable} guna mengakomodasi keterbatasan fisik pelari dan kondisi pencahayaan yang minim.

Hasil pengujian membuktikan bahwa sistem yang dikembangkan telah memenuhi seluruh kriteria keberhasilan. \textit{Functionality Testing} menunjukkan tingkat keberhasilan 100\% pada 26 skenario fungsional. Melalui simulasi \textit{End-to-End Testing}, antarmuka terbukti mampu menangani skalabilitas visualisasi data secara mulus untuk beban 300 tim. Lebih lanjut, hasil \textit{User Acceptance Testing} (UAT) dari 35 responden mencapai skor rata-rata 4,56 dari 5,00, yang menegaskan bahwa antarmuka sistem pelacakan ini sangat fungsional, responsif, dan memberikan pengalaman penggunaan yang sangat baik.
\vspace{0.5cm}

Kata kunci: \textit{Ultra-Marathon}, Pelacakan Peserta, \textit{Front-End}, Aplikasi \textit{Mobile}, Aplikasi Web, \textit{Real-Time Tracking}.
\end{spacing}
